\section{Introduzione}

Il presente documento descrive le fasi finali di implementazione e valutazione del progetto Culturia. Data la notevole ampiezza dell'ecosistema e la presenza di due moduli distinti e separati (il portale gestionale Web e l'applicazione Mobile per gli utenti finali), il team ha optato per una suddivisione mirata dei compiti:
\begin{itemize}
    \item \textbf{Antonio Walter De Fusco} si è occupato dello sviluppo completo della parte Web funzionante. Questo ha richiesto l'implementazione del frontend, del backend (con l'integrazione del database Firebase) e del motore di Intelligenza Artificiale (AI Narrator) per l'estrazione e generazione di contenuti culturali. Si fa presente che questo complesso modulo architetturale è attualmente ancora in fase di sviluppo attivo per scopi di Tesi e di prodotto reale.
    \item \textbf{Gabriele di Palma} si è occupato di realizzare e raffinare un primo prototipo interattivo (mockup basato su tecnologie web) relativo all'applicazione Mobile.
\end{itemize}

Questa convergenza di sforzi ci ha permesso di ottenere entrambi i frontend (Web gestionale e App mobile) a un livello di fedeltà sufficiente per condurre una proficua ed estesa campagna di test di usabilità (Usability Testing) coinvolgendo un panel di utenti reali.
